\documentclass[conference]{IEEEtran}
\IEEEoverridecommandlockouts

\usepackage{cite}
\usepackage{amsmath}
\usepackage{amssymb}
\usepackage{graphicx}
\usepackage{booktabs}
\usepackage{microtype}
\usepackage{hyperref}

\begin{document}

\title{Working title}

\author{\IEEEauthorblockN{Author One}
\IEEEauthorblockA{\textit{Department} \\
\textit{Institution}\\
City, Country \\
author@example.org}
}

\maketitle

\begin{abstract}
One paragraph. Problem, approach, headline result with a number, and the
consequence. Around 150 to 250 words for most IEEE conferences, but check the
call for the exact limit.
\end{abstract}

\begin{IEEEkeywords}
keyword one, keyword two, keyword three
\end{IEEEkeywords}

\section{Introduction}
\label{sec:introduction}

The question, why it is open, and the contributions of this paper.

\section{Related Work}
\label{sec:related}

The neighbouring lines of work and the gap this paper fills.

\section{Method}
\label{sec:method}

Setup, notation, procedure.

\begin{equation}
\label{eq:objective}
\mathcal{L}(\theta) = \frac{1}{n} \sum_{i=1}^{n} \ell\bigl(f_\theta(x_i), y_i\bigr)
\end{equation}

\section{Experiments}
\label{sec:experiments}

Data, baselines, protocol, metrics, and uncertainty.

\begin{table}[t]
  \centering
  \caption{Results on the evaluation suite. Replace with real numbers.}
  \label{tab:results}
  \begin{tabular}{lcc}
    \toprule
    Method & Accuracy (\%) & Latency (ms) \\
    \midrule
    Baseline & 00.0 & 000 \\
    This work & 00.0 & 000 \\
    \bottomrule
  \end{tabular}
\end{table}

\section{Conclusion}
\label{sec:conclusion}

What changed, and the limitation a reader should carry away.

\end{document}
